\documentclass[11pt]{amsart}

\usepackage[T1]{fontenc}
\usepackage{lmodern}
\usepackage{microtype}
\usepackage{mathtools,amssymb,amsthm}
\usepackage{array,booktabs}
\usepackage[hidelinks]{hyperref}

\hypersetup{
  pdftitle={Girth-12 lifts of the voltage graph G_12 for m = 6, 7, 8,
            and nonexistence on its printed arc set for m = 4, 5}
}

\newtheorem{theorem}{Theorem}
\newtheorem{lemma}[theorem]{Lemma}
\newtheorem{proposition}[theorem]{Proposition}
\theoremstyle{remark}
\newtheorem*{remark}{Remark}

\newcommand{\Zm}{\mathbb{Z}_m}
\newcommand{\Gt}{G_{12}}
\newcommand{\Tt}{T_{12}}

\title[Girth-12 lifts of $G_{12}$ for $m = 6,7,8$]
{Girth-12 Lifts of the Voltage Graph $G_{12}$ for $m = 6, 7, 8$,
 and Nonexistence on Its Printed Arc Set for $m = 4, 5$}

\date{}

\subjclass[2020]{05C35, 05C25, 05C38}
\keywords{semicubic cage, biregular cage, girth, voltage graph, lift,
derived graph, exhaustive search}

\begin{document}

\begin{abstract}
Aguilar, Araujo-Pardo and Berman construct a voltage graph $\Gt$ whose
$\Zm$-lift is a $(\{3,m\};12)$-graph of order $49m+3$ for every $m \ge 9$, and
ask, in the second sub-bullet of their open-questions section, for voltage
assignments producing $(\Gt;m)$-graphs for $4 \le m \le 8$ or a proof that none
exist. For $m = 6, 7, 8$ we exhibit voltage assignments on the arcs of $\Gt$
whose lifts are $(\{3,m\};12)$-graphs of orders $297$, $346$, $395$; that
answers the sub-bullet at those three values of $m$. Separately, with the $24$
arcs of their Table 3 held fixed as printed --- with the evident correction to
the duplicated starting leaf --- and only the $24$ voltages varying, we prove
that no assignment gives girth $12$ at $m = 4$ or at $m = 5$: at $m = 4$ by a
four-variable pigeonhole argument requiring no computation, four specific arc
voltages having to be pairwise distinct and nonzero while $\mathbb{Z}_4$ has
only three nonzero residues; at $m = 5$ by exhaustion of the complete cell
$\mathbb{Z}_5^{24}$. We do not settle whether the sub-bullet asks only about
this fixed arc set or about every girth-$12$ graph obtainable from the same
underlying tree; under the wider reading the cells $m = 4, 5$ remain open, so
what is proved here at $m = 4, 5$ is a fixed-arc statement. No order record and
no improved upper bound is claimed at any $m$: the published bounds dominate all
three witnesses.
\end{abstract}

\maketitle

\section{The question}\label{sec:q}

Write $n(\{r,m\};g)$ for the least order of a graph all of whose degrees lie in
$\{r,m\}$, in which both degrees occur, and whose girth is $g$; such a graph is
a \emph{$(\{r,m\};g)$-graph}. Aguilar, Araujo-Pardo and
Berman~\cite{AAB} study these for $r=3$ (\emph{semicubic}) using voltage
graphs. Throughout their paper the voltage group is cyclic: ``\emph{Throughout
this paper, $\Gamma$ is a cyclic group $\mathbb{Z}_m$ with addition as the group
operation}''~\cite{AABarxiv}.

Their Theorem 5.5 reads, verbatim and identically in the
preprint~\cite{AABarxiv} and in the version of record~\cite[p.~25]{AAB}:

\begin{quote}
``The graph ($\Gt$ ; $m$) in $\mathcal{G}_{4t}$ for $t = 3$ given by the voltage
graph $\Gt$ formed by adding the arcs and voltage assignments to $\Tt$ given in
Table 3 has girth 12 for $m \ge 9$. Moreover, it is a $(\{3, m\}; 12)$-graph
with three vertices of degree $m$ and $49m$ vertices of degree 3.''
\end{quote}

The question addressed here is the \textbf{second sub-bullet} under the bullet
``Related to semicubic graphs'' of their Section~6, ``Open questions'', verbatim
from~\cite[p.~29]{AAB} and, identically, from the preprint~\cite{AABarxiv} (the
``$\le$'' is the PDF's glyph and the subscripts are flattened as in the source's
own inline style):

\begin{quote}
``Find voltage assignments to produce ($\Gt$ ; $m$) graphs with the
construction given in Theorem 5.5 for any $4 \le m \le 8$, or show no such
graphs exist.''
\end{quote}

By Theorem 5.5 the object sought for a given $m$ has order $49m+3$, that is
$199, 248, 297, 346, 395$ for $m = 4, \dots, 8$, with three vertices of degree
$m$ and $49m$ of degree $3$, and is the $\Zm$-lift of $\Gt$ under \emph{some}
voltage assignment on the $24$ arcs of Table 3.

\begin{theorem}\label{thm:main}
Let $\Gt$ be the graph of Section~\textup{\ref{sec:obj}} --- the tree $\Tt$ of
Table~\textup{\ref{tab:tree}} together with the $24$ arcs of
Table~\textup{\ref{tab:arcs}}. For each $m \in \{6,7,8\}$ there is a $\Zm$
voltage assignment on those $24$ arcs whose lift is a $(\{3,m\};12)$-graph of
order $49m+3$, namely $297$, $346$, $395$; one such assignment is exhibited for
each $m$ in Section~\textup{\ref{sec:yes}}.
\end{theorem}

\begin{theorem}\label{thm:no}
Let the $24$ arcs of Table~\textup{\ref{tab:arcs}} be held fixed as printed, so
that only the $24$ voltages vary. Then no $\Zm$ voltage assignment on those
arcs has a lift of girth $12$ for $m = 4$
\textup{(}Proposition~\textup{\ref{prop:m4}}, a four-line argument with no
computation\textup{)}, nor for $m = 5$
\textup{(}Section~\textup{\ref{sec:m5}}, by exhaustion of all $5^{24}$
assignments\textup{)}.
\end{theorem}

\medskip
\noindent\textbf{Scope of Theorem~\ref{thm:no}.} Theorem~\ref{thm:main} needs no
caveat: the assignments it exhibits live on the printed arcs, which the
sub-bullet permits on any reading, so at $m = 6, 7, 8$ the sub-bullet is
answered outright. Theorem~\ref{thm:no} is narrower --- the $24$ arcs are held
\emph{fixed as printed} --- and \emph{we do not settle that this is the reading
its authors intended}. What can be said for it is this: Theorem 5.5 defines
$\Gt$ as $\Tt$ plus ``\emph{the arcs \textbf{and voltage assignments} \dots
given in Table 3}'', so the arcs are part of the named object, and~\cite{AAB}
poses a re-arcing question in a \emph{separate} sub-bullet, about a different
graph $H'_{12}$. That is an inference about wording and not a proof of intent;
no author clarification was obtained. Accordingly Theorem~\ref{thm:no} must
always be read as ``\emph{no $\Zm$ voltage assignment on the arc set of $\Gt$
gives girth 12}'' at $m = 4, 5$, and never as the unqualified ``no such graphs
exist''. Under the wider reading, in which each $x$-leaf may be re-paired to
some $y$-leaf and some $z$-leaf, the cells $m = 4, 5$ are untouched here and
remain open; see Section~\ref{sec:open}. The group is $\Zm$ and nothing else, by
the sentence quoted above; in particular the $m=4$ cell is $\mathbb{Z}_4^{24}$
and not also $(\mathbb{Z}_2 \times \mathbb{Z}_2)^{24}$.

\medskip
\noindent Theorem 5.5 is not weakened by any of this. With the printed Table-3
voltages the lift has girth $6,6,8,8,8,10$ at $m = 3,\dots,8$ and $12$ at
$m = 9,\dots,14$ (recomputed by the program of Section~\ref{sec:prog}), so its
threshold $m \ge 9$ is exactly attained for the object it names --- a voltage
graph includes its voltages --- and what Theorem~\ref{thm:main} exhibits are
\emph{other} assignments on the same underlying graph.

\medskip
\noindent\textbf{No record is claimed.} The published upper bounds at the three
cells where we exhibit a graph are $n(\{3,6\};12) \le 243$,
$n(\{3,7\};12) \le 334$ and $n(\{3,8\};12) \le 374$, as tabulated by
Goedgebeur, Jooken and Van den Eede~\cite[girth-$12$ rows of its bound
table]{GJV}; the last is Corollary~3 of~\cite{AAB} itself. Our three witnesses
have orders $297$, $346$, $395$ and are dominated in every case. The source's
own introduction says of this family that ``\emph{the order of these graphs is
bigger than the upper bounds from the identifying vertices construction}''; a
reading of the witnesses as improvements would contradict it.

\medskip
\noindent\textbf{Adjacent work.} \cite{GJV} generalises the gluing theorem
of~\cite{AAB}, and its girth-$12$ table gives bounds for exactly the five cells
at issue, the bound $243$ at $m = 6$ among them; but it uses no voltage graphs
or lifts, and it does not address the existence of a girth-$12$ $\Zm$ voltage
assignment on the arc set of $\Gt$, which is the question here. Exoo and
Jajcay~\cite{EJ} give girth criteria for $\Zm$ lifts of voltage graphs, the
closest work methodologically, but treat neither $\Gt$ nor these five cells;
Exoo~\cite{Exoo04} searches voltage assignments for regular cages, a related
technique on different objects. Araujo-Pardo, Ramos-Rivera and
Jajcay~\cite{ARJ} and Araujo-Pardo, Kiss and Sz\H{o}nyi~\cite{AKS} reach
girth-$12$ bipartite biregular cages through generalized hexagons and designs,
again different objects, and settle nothing about $\Gt$.

\section{The graph $\Gt$, written out in full}\label{sec:obj}

The source gives $\Tt$ by prose and the arcs by its Table 3. Both are
transcribed here completely, so that no figure, no external file and no program
is needed to read what follows. Throughout, ``Table 3'' with no further
qualification means the source's table; the tables of the present note are
Tables~\ref{tab:tree} and~\ref{tab:arcs}, and Table~\ref{tab:arcs} is the
source's Table 3 with the correction of Remark~\ref{rem:typo} applied.

$\Tt$ has $52$ vertices --- $49$ ordinary ones and three \emph{pinned} vertices
$x^*, y^*, z^*$ of degree $1$ --- and $51$ edges, all carrying voltage $0$. A
pinned vertex is a degree-$1$ base vertex: in the lift a \emph{single} vertex
$x^*$ is joined to all $m$ copies of its unique neighbour, so
$\deg(x^*) = m$ and the voltage on that edge is irrelevant.

\begin{table}[ht]
\caption{The $51$ tree edges of $\Tt$, all of voltage $0$.}\label{tab:tree}
\small
\begin{tabular}{@{}l@{\;}l@{}}
\toprule
$X'_3$ ($26$ vertices, $25$ edges)
 & $x^*x$;\quad $x\,x_0$, $x\,x_1$;\quad
   $x_0x_{00}$, $x_0x_{01}$;\quad $x_1x_{10}$, $x_1x_{11}$;\\
 & $x_{01}x_{010}$, $x_{01}x_{011}$;\quad
   $x_{10}x_{100}$, $x_{10}x_{101}$;\quad
   $x_{11}x_{110}$, $x_{11}x_{111}$;\\
 & $x_{010}x_{0100}$, $x_{010}x_{0101}$;\quad
   $x_{011}x_{0110}$, $x_{011}x_{0111}$;\\
 & $x_{100}x_{1000}$, $x_{100}x_{1001}$;\quad
   $x_{101}x_{1010}$, $x_{101}x_{1011}$;\\
 & $x_{110}x_{1100}$, $x_{110}x_{1101}$;\quad
   $x_{111}x_{1110}$, $x_{111}x_{1111}$\\
\midrule
$Y_2$ ($13$ vertices, $12$ edges)
 & $y^*y$;\quad $y\,y_0$, $y\,y_1$;\quad $y_0y_{01}$;\quad
   $y_1y_{10}$, $y_1y_{11}$;\\
 & $y_{01}y_{010}$, $y_{01}y_{011}$;\quad
   $y_{10}y_{100}$, $y_{10}y_{101}$;\quad
   $y_{11}y_{110}$, $y_{11}y_{111}$\\
\midrule
$Z_2$ ($13$ vertices, $12$ edges)
 & the same with $z$ in place of $y$\\
\midrule
join edges & $x_{00}y_0$, \quad $x_{00}z_0$\\
\bottomrule
\end{tabular}
\end{table}

Equivalently, and this is how~\cite{AAB} defines it: the vertices of $X'_3$ are
the binary strings of length at most $4$ with every string beginning $00$
deleted except $00$ itself (so $x_{00}$ is a leaf of the tree part), plus the
pinned $x^*$ above the root; the vertices of $Y_2 = Z_2$ are the binary strings
of length at most $3$ with every string beginning $00$ deleted, plus $y^*$;
and $x_{00}$ is joined to $y_0$ and to $z_0$. The twelve depth-$5$ leaves
$x_{0100}, \dots, x_{1111}$ are exactly the twelve ``Starting leaf'' entries of
Table 3, and the six leaves of each of $Y_2$, $Z_2$ are its ``End'' entries.

\begin{table}[ht]
\caption{The $24$ arcs, with the authors' own symbol names and the printed
Table-3 voltages. Reading the symbols in the order
$a,b,c,\dots,\varepsilon$ reproduces both printed voltage rows.}\label{tab:arcs}
\small
\begin{tabular}{@{}clr@{\qquad}clr@{}}
\toprule
sym & arc & volt. & sym & arc & volt.\\
\midrule
$a$ & $x_{0100}\to y_{010}$ & $+1$ & $b$ & $x_{0100}\to z_{110}$ & $+2$\\
$c$ & $x_{0101}\to y_{011}$ & $+2$ & $d$ & $x_{0101}\to z_{111}$ & $+1$\\
$e$ & $x_{0110}\to y_{100}$ & $+1$ & $f$ & $x_{0110}\to z_{100}$ & $-1$\\
$g$ & $x_{0111}\to y_{101}$ & $+2$ & $h$ & $x_{0111}\to z_{101}$ & $+3$\\
$j$ & $x_{1000}\to y_{010}$ & $+2$ & $k$ & $x_{1000}\to z_{010}$ & $-1$\\
$l$ & $x_{1001}\to y_{011}$ & $+1$ & $p$ & $x_{1001}\to z_{011}$ & $+3$\\
$q$ & $x_{1010}\to y_{110}$ & $+1$ & $r$ & $x_{1010}\to z_{100}$ & $+1$\\
$s$ & $x_{1011}\to y_{111}$ & $+2$ & $t$ & $x_{1011}\to z_{101}$ & $-1$\\
$u$ & $x_{1100}\to y_{100}$ & $+2$ & $v$ & $x_{1100}\to z_{010}$ & $-3$\\
$w$ & $x_{1101}\to y_{101}$ & $+1$ & $\alpha$ & $x_{1101}\to z_{011}$ & $-2$\\
$\beta$ & $x_{1110}\to y_{110}$ & $+3$ & $\gamma$ & $x_{1110}\to z_{110}$ & $-1$\\
$\delta$ & $x_{1111}\to y_{111}$ & $-1$ & $\varepsilon$ & $x_{1111}\to z_{111}$ & $+2$\\
\bottomrule
\end{tabular}
\end{table}

The symbol names $a, b, \dots, \varepsilon$ (the list skips $i, m, n, o$) are
the authors' own, taken from a commented-out substitution list in the
preprint's source. Reading the $24$ symbols in the
interleaved order above reproduces both printed voltage rows exactly ---
$x\!\to\!y$: $1,2,1,2,2,1,1,2,2,1,3,-1$ and $x\!\to\!z$:
$2,1,-1,3,-1,3,1,-1,-3,-2,-1,2$ --- which is an internal checksum on the table.

\begin{remark}[a forced correction to the published Table 3]\label{rem:typo}
The \textbf{eighth} ``Starting leaf'' entry is printed $x_{1001}$ in both
sub-rows, duplicating the sixth; it must read $x_{1011}$, as in
Table~\ref{tab:arcs}. This is a defect of the published paper, present in the
preprint \texttt{.tex} and on p.~26 of the version of record alike. Three
independent confirmations: $x_{1011}$ is the only one of the twelve $4$-bit
$x$-leaf labels otherwise missing; the paper's own commented-out earlier
version of the same assignment, in the preprint's source, reads
``$(x_{1011}y_{111}) = 3$, $(x_{1011}z_{101}) = 9$'', and $y_{111}$, $z_{101}$
are exactly the printed eighth ``End'' entries; and under the literal printed
reading $x_{1001}$ would have degree $5$ and $x_{1011}$ degree $1$, so the
object would not be the one Theorem 5.5 describes, whereas with the correction
every non-pinned vertex has degree $3$ and every $y$- and $z$-leaf receives
exactly two arcs. All numbers below use the corrected reading.
\end{remark}

With Table~\ref{tab:arcs} in place, $\Gt$ has $|V| = 52$, $|E| = 75$, cycle rank
$\beta = 75 - 52 + 1 = 24$, degree multiset $\{1^3, 3^{49}\}$ with the three
degree-$1$ vertices exactly $x^*, y^*, z^*$; it is connected, bipartite, and has
no $4$-cycles. All three pinned-to-pinned tree distances equal $6$, so every
lift already contains a $12$-cycle: ``girth $12$'' below means exactly ``no
cycle shorter than $12$'', never ``girth at least $12$''.

\section{The decision criterion}\label{sec:crit}

Let $v$ assign an element of $\Zm$ to each of the $24$ arcs, tree edges being
$0$, and let $\mathcal{C}$ be the set of simple cycles of $\Gt$ of length at
most $10$. Traversing a cycle, an arc contributes $+v$ if crossed in its
direction and $-v$ otherwise.

\begin{lemma}\label{lem:crit}
The lift $(\Gt;m)_v$ has girth $12$ if and only if
$\sum_{e \in C} \pm v(e) \not\equiv 0 \pmod m$ for every $C \in \mathcal{C}$.
\end{lemma}

\begin{proof}
Necessity is elementary: if a simple cycle $C$ of length $L \le 10$ has
$\sum_C v \equiv 0$, then $C$ lifts to a simple $L$-cycle, and $L < 12$.

For sufficiency, three families of base walks could lift to a short cycle:
simple cycles, which the condition handles; lollipop walks through a pinned
vertex; and non-reversing non-cycle walks. The proof of Theorem 5.5
in~\cite[p.~25]{AAB} disposes of the last two, verbatim: ``\emph{By
construction of $\Tt$, all lollipop walks lift to long cycles. Next, we need to
consider non-reversing non-cycle walks in the voltage graph that may lift to
short cycles. From Observation \textup{[nonCycleWalks]}, any such walk must
involve at least one 4-cycle. However, analysis of the cycles in $\Gt$ shows
that, in fact, there are no 4-cycles (by construction), so there are no short
walks of this sort that lift to short cycles.}'' Both disposals are properties
of the \emph{arc set alone}, so they hold for every assignment on those arcs and
not only for the printed one. Directly: a pinned-vertex lollipop with tree-stem length $d$ and loop length $c$ lifts to a
cycle of length $2d + c + 2 \ge 12$, the two minima being $(d,c) = (1,8)$ via
the $8$-cycle
$y_0\,y_{01}\,y_{010}\,x_{0100}\,x_{010}\,x_{01}\,x_0\,x_{00}\,y_0$ and
$(d,c) = (2,6)$, each giving exactly $12$; and $\Gt$ has no $4$-cycle.
\end{proof}

Only necessity is used for the nonexistence proofs, so exhausting the assignment
space against the condition and finding nothing \emph{proves} nonexistence:
dropping a walk family can only weaken the constraint set, which is safe for a
NO. Theorem~\ref{thm:main} does not rely on Lemma~\ref{lem:crit} at all --- every
witness in Section~\ref{sec:yes} is verified by building the full lift and
computing its girth.

Two facts about $\mathcal{C}$ and about the parametrisation:

\begin{proposition}\label{prop:census}
$|\mathcal{C}| = 126$, with $12$ cycles of length $6$, $24$ of length $8$ and
$90$ of length $10$; counted with both orientations this is $252$, the number
published in the proof of Theorem 5.5 \textup{(}``\emph{There are 252 short
cycles that are not formed by going back and forth along a single
edge}''\textup{)}.
\end{proposition}

\begin{proposition}\label{prop:gauge}
The $24$ arc voltages are a complete parametrisation of the $\Zm$ lifts of
$\Gt$ up to isomorphism. Consequently the $m = 4$ and $m = 5$ cells are exactly
$\mathbb{Z}_4^{24}$ and $\mathbb{Z}_5^{24}$.
\end{proposition}

\begin{proof}
This is the standard spanning-tree normalisation for voltage graphs (Gross and
Tucker), applied here. A pinned vertex is joined to all $m$ copies of its
neighbour, so its edge voltage is irrelevant: that fixes $3$ of the $51$ tree
edges. The gauge move $v(e) \mapsto v(e) + f(\mathrm{head}) - f(\mathrm{tail})$,
for $f : V_{\mathrm{ord}} \to \Zm$, induces the lift isomorphism
$w^i \mapsto w^{i + f(w)}$; the $49$ ordinary vertices span a tree with $48$
edges, and choosing $f$ along that tree sets all $48$ of those voltages to $0$.
Since $48 + 3 = 51 = |E(\Tt)|$, all tree voltages may be taken $0$ and the
remaining $\beta = 24$ arc voltages are free.
\end{proof}

Of this proposition the program of Section~\ref{sec:prog} checks only the counts
$49$, $48$ and $48+3 = 51 = |E(\Tt)|$; the normalisation argument itself is the
standard one just cited and is not re-verified there. Nothing in
Theorem~\ref{thm:no} depends on the ``up to isomorphism'' clause: the two
nonexistence proofs quantify over \emph{all} of $\mathbb{Z}_4^{24}$ and
$\mathbb{Z}_5^{24}$ with no quotient taken.

Exactly $14$ of the $24$ arcs are forced nonzero, namely those whose
fundamental cycle has length at most $10$: $a, b, c, d, e, f, g, h, j, k, l, p,
v, \alpha$, with $a$ and $c$ at length $8$ and the other twelve at length $10$.
The remaining ten, $q, r, s, t, u, w, \beta, \gamma, \delta, \varepsilon$, have
fundamental cycles of length $12$ and may legitimately be $0$; zero entries do
occur in the witnesses below.

\section{Existence at $m = 6, 7, 8$}\label{sec:yes}

Each witness is $24$ residues, listed in the symbol order of
Table~\ref{tab:arcs}. The base graph is Section~\ref{sec:obj} and the lift is
the standard $\Zm$ derived graph, so a reader needs nothing further.

\medskip
\noindent\textbf{$m = 6$, a $(\{3,6\};12)$-graph of order $297$.}
\[
\begin{aligned}
&a{=}1,\ b{=}3,\ c{=}3,\ d{=}4,\ e{=}2,\ f{=}1,\ g{=}4,\ h{=}2,
 \ j{=}2,\ k{=}3,\ l{=}1,\ p{=}5,\\
&q{=}3,\ r{=}3,\ s{=}2,\ t{=}5,\ u{=}5,\ v{=}1,\ w{=}3,\ \alpha{=}2,
 \ \beta{=}0,\ \gamma{=}5,\ \delta{=}1,\ \varepsilon{=}3.
\end{aligned}
\]

\noindent\textbf{$m = 7$, a $(\{3,7\};12)$-graph of order $346$.}
\[
\begin{aligned}
&a{=}1,\ b{=}3,\ c{=}3,\ d{=}4,\ e{=}4,\ f{=}1,\ g{=}6,\ h{=}2,
 \ j{=}2,\ k{=}3,\ l{=}1,\ p{=}4,\\
&q{=}2,\ r{=}0,\ s{=}4,\ t{=}4,\ u{=}2,\ v{=}1,\ w{=}0,\ \alpha{=}5,
 \ \beta{=}3,\ \gamma{=}0,\ \delta{=}6,\ \varepsilon{=}5.
\end{aligned}
\]

\noindent\textbf{$m = 8$, a $(\{3,8\};12)$-graph of order $395$}: the $m=7$
vector read in $\mathbb{Z}_8$, with $g = 7$ in place of $6$.

\medskip
In each case the lift has $49m+3$ vertices and $75m$ edges, three vertices of
degree $m$ and $49m$ of degree $3$, and girth exactly $12$. The three
occurrences of $0$ at $m=7,8$ are legal: $r$, $w$, $\gamma$ are among the ten
arcs whose fundamental cycles have length $12$.

\section{Nonexistence at $m = 4$: a pigeonhole, with no computation}
\label{sec:m4}

\begin{proposition}\label{prop:m4}
No $\mathbb{Z}_4$ voltage assignment on the $24$ arcs of $\Gt$ has a lift of
girth $12$. The same argument re-proves the case $m = 3$.
\end{proposition}

\begin{proof}
Consider the four arcs into $z_{010}$ and $z_{011}$, in the notation of
Table~\ref{tab:arcs}:
\begin{align*}
k &= v(x_{1000}\!\to\! z_{010}), & v &= v(x_{1100}\!\to\! z_{010}),\\
p &= v(x_{1001}\!\to\! z_{011}), & \alpha &= v(x_{1101}\!\to\! z_{011}).
\end{align*}

\emph{(i) Each is nonzero.} In $\Tt$ we have $d(x_{1000}, z_{010}) = 9$: six
steps up to $x_{00}$, then three down, along
\[
x_{1000}\ x_{100}\ x_{10}\ x_1\ x\ x_0\ x_{00}\ z_0\ z_{01}\ z_{010}.
\]
So the arc $k$ closes a cycle of length $1 + 9 = 10 \le 10$ whose only arc is
$k$, namely
\[
x\ x_0\ x_{00}\ z_0\ z_{01}\ z_{010}\ x_{1000}\ x_{100}\ x_{10}\ x_1\ x .
\]
By Lemma~\ref{lem:crit} (necessity), $k \not\equiv 0$. The three other
distances are $9$ as well, so $p, v, \alpha \not\equiv 0$.

\emph{(ii) They are pairwise distinct.} Each of the six pairs is joined by a
cycle of length at most $10$ traversing one of the two arcs forwards and the
other backwards, using $d(x_{1000},x_{1001}) = d(x_{1100},x_{1101}) = 2$, the
four cross distances $= 6$, and $d(z_{010},z_{011}) = 2$:
\begin{center}\small
\begin{tabular}{@{}lll@{}}
\toprule
pair & length & cycle\\
\midrule
$k,p$ & $6$ & $x_{100}\,x_{1000}\,z_{010}\,z_{01}\,z_{011}\,x_{1001}$\\
$v,\alpha$ & $6$ & $x_{110}\,x_{1100}\,z_{010}\,z_{01}\,z_{011}\,x_{1101}$\\
$k,v$ & $8$ & $x_1\,x_{10}\,x_{100}\,x_{1000}\,z_{010}\,x_{1100}\,x_{110}\,x_{11}$\\
$p,\alpha$ & $8$ & $x_1\,x_{10}\,x_{100}\,x_{1001}\,z_{011}\,x_{1101}\,x_{110}\,x_{11}$\\
$k,\alpha$ & $10$ & $x_1\,x_{10}\,x_{100}\,x_{1000}\,z_{010}\,z_{01}\,z_{011}\,x_{1101}\,x_{110}\,x_{11}$\\
$p,v$ & $10$ & $x_1\,x_{10}\,x_{100}\,x_{1001}\,z_{011}\,z_{01}\,z_{010}\,x_{1100}\,x_{110}\,x_{11}$\\
\bottomrule
\end{tabular}
\end{center}
Each such cycle has voltage sum $\pm(v_1 - v_2)$ for the two arcs concerned, so
necessity gives $k \ne p$, $k \ne v$, $k \ne \alpha$, $p \ne v$,
$p \ne \alpha$, $v \ne \alpha$ in $\Zm$.

\emph{(iii)} Hence $k, p, v, \alpha$ are four pairwise distinct nonzero
elements of $\Zm$, which requires $m - 1 \ge 4$. But $\mathbb{Z}_4$ has only
three nonzero residues, and $\mathbb{Z}_3$ only two. Pigeonhole.
\end{proof}

\section{Nonexistence at $m = 5$}\label{sec:m5}

\begin{proposition}\label{prop:m5}
No $\mathbb{Z}_5$ voltage assignment on the $24$ arcs of $\Gt$ has a lift of
girth $12$.
\end{proposition}

The cell is $\mathbb{Z}_5^{24}$ by Proposition~\ref{prop:gauge}, that is
$5^{24} = 59{,}604{,}644{,}775{,}390{,}625$ points. Branch-and-prune against the
$126$ necessary conditions of Lemma~\ref{lem:crit}, over the \emph{complete}
space with no gauge quotient and no symmetry break of any kind, exhausts it and
returns no solution. Only necessity is used, so the negative cannot be an
artefact of an over-strong constraint set. This case is machine-only here: we
know of no short hand argument for it.

\section{The accompanying program}\label{sec:prog}

This folder contains \texttt{verify.py} and \texttt{verify.output.txt}, a
recorded run of it: Python~3.9, standard library only, no solver and no external
data file, $80$ checks, all passing. It is built from Section~\ref{sec:obj}
alone and does not parse this note. The quantities transcribed out of the note
and re-derived there are: the tree by two independent constructions, and the
invariants of Section~\ref{sec:obj};
the census of Proposition~\ref{prop:census} by two algorithms (path-DFS with
min-index canonicalisation, and XOR of fundamental cycles over the co-tree
arcs), agreeing on edge sets; the impossibility of the literal printed Table 3
(Remark~\ref{rem:typo}); published invariants of~\cite{AAB} --- the
``$252$ short cycles'' as $126 \times 2$, the absolute voltage sums lying in
``$\{1,2,\dots,8\}$'', the absence of $4$-cycles, the order $49m+3$ with three
vertices of degree $m$, and the girth profile of the printed table reported in
Section~\ref{sec:q}; the girth of each
witness of Section~\ref{sec:yes} on the full lift, by two algorithms
(BFS from every vertex, and deleting each edge and BFS-ing between its
endpoints); the ten constraints of Proposition~\ref{prop:m4}, each identified
with one of the $126$ enumerated cycle forms; and the exhaustions of
$\mathbb{Z}_4^{24}$ and
$\mathbb{Z}_5^{24}$, all $24$ variables branched, all $126$ forms enforced, no
gauge quotient and no symmetry break.

What it does not do. It does not re-verify the normalisation argument of
Proposition~\ref{prop:gauge}, only its counts; it says nothing about the wide
variant of Section~\ref{sec:open}; and it makes no statement about published
order records, of which this note claims none. The recorded run also prints
checks on matters this note makes no claim about, and it makes no claim to have
enumerated everything stated here.

\section{What is not settled}\label{sec:open}

\begin{enumerate}
\item \textbf{The wide reading is open at $m = 4$ and $m = 5$, and which reading
the sub-bullet intends is not settled here.}
Whether a \emph{different} pairing of the twelve $x$-leaves to
the $y$- and $z$-leaves --- one of each, no $4$-cycles --- admits a girth-$12$
$\mathbb{Z}_4$ or $\mathbb{Z}_5$ assignment is untouched. In favour of the narrow
reading, the sibling sub-bullets of~\cite[p.~28--29]{AAB} spell the wide variant
out explicitly (``\emph{either a different voltage assignment for the same arcs
as in $H_{10}$, or a different assignment of arcs and voltages using the same
`underlying' tree}'') and ours contains no such device; but that is an inference
about wording, not an author statement, so a reader who prefers the wide reading
should treat Theorem~\ref{thm:no} as a statement about the $\Gt$ arc set only,
and the wide question at $m = 4, 5$ as open. The issue does not arise at
$m = 6, 7, 8$, where Theorem~\ref{thm:main} answers the sub-bullet on either
reading.
\item \textbf{Nothing is claimed about $m \ge 9$} beyond reproducing Theorem 5.5,
nor about the sibling $H_{10}$ and $H_{12}$ sub-bullets of the same
open-questions section, nor about the $(3,14)$-cage bullet, nor about any order
record or upper bound at any $m$.
\item \textbf{The Table-3 correction is ours.} Remark~\ref{rem:typo} is a
defect of the published paper; the correction is forced three ways, but it is
a correction and a referee should see it as such.
\end{enumerate}

\begin{thebibliography}{9}

\bibitem{AAB}
F.~Aguilar, G.~Araujo-Pardo, and L.~W. Berman,
\emph{Semicubic cages and small graphs of even girth derived from voltage
graphs},
Ars Math. Contemp. \textbf{26} (2026), \#P3.03.
\href{https://doi.org/10.26493/1855-3974.3116.35e}{doi:10.26493/1855-3974.3116.35e}.
Received 5 May 2023, accepted 6 August 2025, published online 2 June 2026.
Page references are to the journal galley.

\bibitem{AABarxiv}
F.~Aguilar, G.~Araujo-Pardo, and L.~W. Berman,
\emph{Semicubic cages and small graphs of even girth from voltage graphs},
\href{https://arxiv.org/abs/2305.03290}{arXiv:2305.03290v1}, 5 May 2023.
Quotations of commented-out material are from its single e-print source file
\texttt{AABsemicubiccagesArxive.tex}. The version of record~\cite{AAB} carries
``derived'' in its title; this version does not.

\bibitem{GJV}
J.~Goedgebeur, J.~Jooken, and T.~Van den Eede,
\emph{Computational methods for finding bi-regular cages},
\href{https://arxiv.org/abs/2411.17351}{arXiv:2411.17351v1}, 26 November 2024.
The girth-$12$ rows of its bound table are the published landscape for the cells
at issue here; the bound $374$ at $m = 8$ is attributed there to Corollary~3
of~\cite{AAB}. The bounds are quoted at second hand through that table.

\bibitem{EJ}
G.~Exoo and R.~Jajcay,
\emph{On the girth of voltage graph lifts},
European J. Combin. \textbf{32} (2011).
\href{https://doi.org/10.1016/j.ejc.2010.12.003}{doi:10.1016/j.ejc.2010.12.003}.

\bibitem{Exoo04}
G.~Exoo,
\emph{Voltage graphs, group presentations and cages},
Electron. J. Combin. \textbf{11} (2004).
\href{https://doi.org/10.37236/1843}{doi:10.37236/1843}.

\bibitem{ARJ}
G.~Araujo-Pardo, A.~Ramos-Rivera, and R.~Jajcay, on bipartite biregular cages,
\href{https://arxiv.org/abs/1907.11568}{arXiv:1907.11568}. Cited by e-print
number and authors only; the title is not quoted here.

\bibitem{AKS}
G.~Araujo-Pardo, Gy.~Kiss, and T.~Sz\H{o}nyi, on bipartite biregular cages via
generalized polygons and designs,
\href{https://arxiv.org/abs/2310.12137}{arXiv:2310.12137}. Cited by e-print
number and authors only; the title is not quoted here.

\end{thebibliography}

\end{document}
